\chapter{Two Oracles Play Rock--Paper--Scissors}
\label{chap:two-oracles-rps}

\begin{quote}
\textbf{Author's Intuition.}

\emph{This section is on how do two people that can see eachother moves in advance can cooperate.}
\end{quote}

The sentence above is preserved exactly as the original research question. In
this chapter we study one narrow interpretation of it. We do not assume literal
prophecy. Instead, ``seeing a move in advance'' means learning the other
participant's plaintext choice before fixing one's own. The purpose of the
protocol is to remove that timing advantage while preserving each
participant's move until the reveal phase. Whether the oracle metaphor also
contains a deeper strategic problem remains open.

\begin{abstract}
This chapter gives an educational two-party commit--reveal protocol for
Rock--Paper--Scissors. Alice and Bob each fix a private move, exchange
session-bound commitments, and reveal only after both commitments have been
accepted. A valid reveal is checked against the corresponding prior
commitment. The construction replaces a trusted referee's early collection of
plaintext moves, but it does not reproduce a referee's ability to release both
moves fairly. In particular, the participant who receives the first reveal can
learn the result and then refuse to reveal. The chapter therefore treats
selective abort as a central counterexample, not as a minor implementation
detail. The companion C code is an instructional state machine, not a
production protocol, and the exact hash-based commitment has a limited
random-oracle-style security argument rather than a formal proof.
\end{abstract}

\section{Cooperation Problem}

Rock--Paper--Scissors is normally described as simultaneous choice. Over a
network, however, messages arrive in an order. If Alice sends
\texttt{rock} in plaintext and Bob receives it before choosing, Bob can choose
\texttt{paper}. Reversing the order only reverses the advantage.

Alice and Bob have a shared interest in obtaining a well-defined game, but
their interests diverge over its outcome. Neither should have to reveal first
without protection, and neither should have to trust the other to ignore an
early plaintext move. The protocol question is therefore:

\begin{quote}
Can both participants fix their moves before either learns the other's move,
and can each later verify that an accepted reveal matches the earlier choice?
\end{quote}

This is a chronology problem. It is not a method for choosing a wise move,
making either participant honest, or making completion rational.

\section{Trusted-Referee and Plaintext Baselines}

With a trusted referee, Alice and Bob privately submit their moves. The referee
waits for both, fixes them, and releases both moves or the result. The referee
provides three services:

\begin{enumerate}
    \item it receives each move without disclosing it to the other participant;
    \item it prevents either move from changing after submission; and
    \item it can release the result to both participants together.
\end{enumerate}

Plaintext exchange provides none of these services without an ordering
advantage. Commit--reveal protocols have a long history in remote joint-choice
problems; Blum's coin-flipping work is an early point of reference
\cite{blum1983coinflipping}. That citation records prior context. It is not a
claim that the construction in this chapter is Blum's exact protocol.

The present protocol replaces the referee's first two services under explicit
computational assumptions. It does not replace the third. That difference is
the boundary between simultaneous choice and fair output.

\section{Parties, Moves, Secrets, and Outputs}

There are exactly two fixed roles:

\begin{itemize}
    \item Alice, written \(A\), with private move \(m_A\);
    \item Bob, written \(B\), with private move \(m_B\).
\end{itemize}

The canonical move encoding is

\[
    \mathsf{Rock}=0,\qquad
    \mathsf{Paper}=1,\qquad
    \mathsf{Scissors}=2.
\]

Before the round begins, both parties agree on the protocol version, the
commitment version, one public 32-byte session identifier, their role
assignment, a nonzero 32-bit round number, the move encoding, and their local
timeout policies.

Before reveal, each participant's move and fresh 32-byte commitment nonce are
private. The protocol version, session identifier, round, roles, commitment
digests, message timing, and message lengths are public. At successful
completion, both moves, both nonces, and the result are public. Repeated games
therefore reveal a history from which strategies might be inferred.

The session identifier separates executions. It is not a secret, a signature,
an identity, or proof of freshness by itself.

For completed moves, define

\[
    d=(m_A-m_B+3)\bmod 3.
\]

The result is a tie when \(d=0\), Alice wins when \(d=1\), and Bob wins when
\(d=2\).

\section{Threat Model, Channels, and Setup}

The intended analysis has a classical, computationally bounded adversary that
statically corrupts at most one participant. A corrupted participant may
choose any legal move or nonce; send malformed, duplicate, replayed,
reordered, or context-mismatched messages; send a random commitment for which
it knows no opening; open inconsistently; delay; remain silent; or abort at any
time.

The transport is assumed to provide origin authentication and message
integrity for messages accepted as coming from the peer. The transport may
still expose, delay, drop, duplicate, or reorder messages. Identity
establishment and authenticated transport are external assumptions: this
chapter does not implement them. In particular, a role byte and a session
identifier do not authenticate a human being.

An honest participant is assumed to have:

\begin{itemize}
    \item a functioning operating-system cryptographic random source;
    \item protected local memory for its move and nonce until reveal;
    \item an agreed and unique protocol context; and
    \item eventual delivery of every valid message when honest completion is
    being claimed.
\end{itemize}

Adaptive corruption, endpoint malware, memory disclosure, coercion, traffic
analysis, crash recovery, concurrent sessions, compromised randomness,
quantum attackers, and side channels beyond digest-comparison discipline are
outside the model.

\section{Idealized Referee and Exact Real Contract}

The idealized referee receives both moves privately and releases the result to
both parties together. It is a comparison baseline, not a functionality that
this chapter claims to realize.

The exact real contract is weaker: \emph{simultaneous choice with abort}.

\begin{enumerate}
    \item Each honest participant fixes its own move and nonce before
    processing a peer message.
    \item An honest participant reveals only after it has sent its own
    commitment and accepted the peer's context-valid commitment.
    \item A peer reveal is accepted only when it opens the previously accepted
    commitment in the same protocol, session, round, and role context.
    \item The participant-local API reports a completed outcome only after the
    local reveal has been emitted and the peer's reveal has validated.
    \item A peer, host, or network may nevertheless prevent completion.
\end{enumerate}

For this chapter, \emph{strong fairness} means that either both parties obtain
the result or neither does. \emph{Guaranteed output delivery} means that an
honest party eventually obtains its output despite peer deviation. The real
protocol provides neither property.

\section{Non-Goals}

This slice does not provide:

\begin{itemize}
    \item networking, identity establishment, PKI, or authenticated transport;
    \item proof at commit time that the sender knows a valid opening;
    \item strong fairness, guaranteed output delivery, or denial-of-service
    resistance;
    \item deposits, penalties, a blockchain, trusted hardware, or a dispute
    authority;
    \item zero knowledge, non-repudiation, deniability, or coercion resistance;
    \item secrecy after reveal or protection against strategy inference across
    repeated games; or
    \item production, post-quantum, adaptive, or composable security.
\end{itemize}

It also does not claim that committing to a move makes that move random,
rational, cooperative, or good.

\section{Commitment Primitive}

Commitment schemes are commonly described through two properties:
\emph{hiding}, which protects the committed value before opening, and
\emph{binding}, which prevents a committer from opening the same commitment as
two different values. Naor gives a formal commitment construction from
pseudorandomness \cite{naor1991bitcommitment}. That work supplies background
for the vocabulary, but it does not prove the different hash-based construction
used here.

For public context \(x\), move byte \(m\), and fresh 32-byte nonce \(r\), the
educational construction computes

\[
    C=\operatorname{BLAKE2b}_{256}
      \bigl(\operatorname{EncodeV1}(x,r,m)\bigr).
\]

The exact 16-byte commitment-domain string is the following ASCII text, without
a terminating NUL:

\begin{center}
\texttt{AC-COMMITMENT-V1}
\end{center}

The 16-byte RPS protocol identifier contains ASCII \texttt{AC-RPS} followed by
ten zero bytes.
The canonical commitment preimage is:

\begin{verbatim}
scheme-domain[16]
scheme-version:u16be
protocol-id-length:u16be
protocol-id[16]
protocol-version:u16be
session-id-length:u16be
session-id[32]
round:u32be
committer-role:u8
recipient-role:u8
payload-type:u16be
nonce-length:u16be
nonce[32]
payload-length:u64be
payload[payload-length]
\end{verbatim}

All integer fields are unsigned and big-endian. Version 1 fixes the scheme,
protocol, session, and nonce lengths but encodes the lengths explicitly so the
byte sequence retains one meaning. No native C structure, padding byte,
pointer, NUL terminator, or ambiguous text concatenation is hashed.

The implementation uses a 32-byte unkeyed BLAKE2b digest. RFC 7693 specifies
BLAKE2b and labels this parameter set with a \(2^{128}\) collision-security
target \cite{saarinen2015blake2}. The RFC also declines to make an independent
security assertion about the algorithm. The parameter label is therefore not
a proof of this commitment construction. A fresh nonce matters because the
move has only three possible values. A bare digest of the move would allow
anyone to try all three possibilities.

\subsection{A note on references and dependencies}

A DOI in the bibliography is a persistent identifier for a paper. An RFC is a
technical specification. Neither is code that the program imports. The
protocol definition above is independent of any library. The C companion uses
only one external cryptographic library, libsodium, which was already a project
dependency. Its documentation identifies the generic-hash backend as BLAKE2b
and documents its random-byte interface
\cite{libsodium-generichash,libsodium-randombytes}. Protocol state, message
encoding, and game logic remain ordinary C11.

\section{Commit--Commit--Reveal--Reveal}

Each participant maintains a separate local state. No protocol object contains
both private moves.

\begin{enumerate}
    \item Alice and Bob independently initialize the same public context. Each
    fixes its local move and samples a fresh secret nonce before processing a
    peer message.
    \item Alice emits \(C_A\), and Bob emits \(C_B\). Either commitment may
    arrive first.
    \item Each side parses the peer commitment and checks protocol version,
    session, round, sender, recipient, type, and length before accepting it.
    \item Only after both commitments exist may an honest participant emit a
    reveal containing its move and nonce.
    \item On receiving a reveal, the participant reconstructs the exact
    commitment context and recomputes the digest.
    \item A mismatched move, nonce, digest, role, session, round, or version is
    rejected. A valid opening records the peer move.
    \item Once the local reveal has been emitted and the peer opening has
    validated, the participant computes the outcome.
    \item A local abort, accepted peer abort, timeout, or established in-session
    violation terminates the local execution without a completed outcome.
\end{enumerate}

Accepting \(C_A\) or \(C_B\) does not prove that its sender knows a valid
opening. Any 32-byte string has valid commitment-message syntax. Knowledge is
tested only when an opening arrives.

\section{Messages, State Machine, and Transcript}

Every message has this canonical envelope:

\begin{verbatim}
magic[4]                 = "ACRP"
wire-version:u16be       = 1
session-id[32]
round:u32be
sender-role:u8           = 1 (Alice) or 2 (Bob)
recipient-role:u8        = 1 (Alice) or 2 (Bob)
message-type:u8
payload-length:u16be
payload[payload-length]
\end{verbatim}

The header is exactly 47 bytes. Version 1 has three message types:

\begin{itemize}
    \item \texttt{COMMIT} (type 1): a 32-byte commitment digest, 79 bytes
    total;
    \item \texttt{REVEAL} (type 2): one move byte followed by a 32-byte
    nonce, 80 bytes total;
    \item \texttt{ABORT} (type 3): one enumerated reason byte, 48 bytes
    total.
\end{itemize}

Version 1 names reason values \texttt{UNSPECIFIED}=1,
\texttt{APPLICATION\_REQUEST}=2, and \texttt{PROTOCOL\_POLICY}=3. A reason is
the sender's declaration; it is not independently verified evidence of cause
or motive.

The parser rejects truncation, trailing bytes, mismatched lengths, unknown
values, equal sender and recipient roles, wrong recipient, zero round, and an
all-zero session identifier. It does not cast a byte buffer to a C structure.

The participant-local states are

\begin{quote}
\raggedright
\texttt{READY}, \texttt{COMMIT\_SENT}, and
\texttt{COMMIT\_RECEIVED};\\
\texttt{BOTH\_COMMITTED}, \texttt{REVEAL\_SENT}, and
\texttt{REVEAL\_RECEIVED};\\
\texttt{COMPLETE}, \texttt{ABORTED\_LOCAL}, and
\texttt{ABORTED\_PEER};\\
\texttt{TIMED\_OUT} and \texttt{FAILED}.
\end{quote}

The central transitions are:

\begin{center}
\small
\begin{tabular}{|>{\raggedright\arraybackslash}p{0.22\textwidth}|
                  >{\raggedright\arraybackslash}p{0.42\textwidth}|
                  >{\raggedright\arraybackslash}p{0.22\textwidth}|}
\hline
State & Accepted action or peer message & Next state \\
\hline
\texttt{READY} & send local commitment & \texttt{COMMIT\_SENT} \\
\hline
\texttt{READY} & receive peer commitment & \texttt{COMMIT\_RECEIVED} \\
\hline
one commitment present & add the other valid commitment & \texttt{BOTH\_COMMITTED} \\
\hline
\texttt{BOTH\_COMMITTED} & send local reveal & \texttt{REVEAL\_SENT} \\
\hline
\texttt{BOTH\_COMMITTED} & receive valid peer reveal & \texttt{REVEAL\_RECEIVED} \\
\hline
one reveal present & add the other valid reveal & \texttt{COMPLETE} \\
\hline
any nonterminal state & local abort, peer abort, or local timeout & matching terminal state \\
\hline
\end{tabular}
\end{center}

Version 1 fails closed on a syntactically valid in-session premature reveal,
duplicate, replay, conflicting commitment, or invalid opening. It does not
buffer and recover from arbitrary network reordering. This protects state
meaning at the cost of availability.

The canonical transcript stores accepted raw messages in fixed logical slots:
Alice commitment, Bob commitment, Alice reveal, Bob reveal, Alice abort, and
Bob abort. Its header contains magic \texttt{ACRT}, transcript version,
session identifier, round, and a six-slot count. Each fixed-order slot contains
a slot identifier, a presence byte, a 16-bit big-endian message length, and the
canonical message when present. An absent slot has presence zero and length
zero. Six slot headers plus the largest possible messages give a version-1
maximum of 481 bytes. It does not store rejected packets or local arrival
order. Honest completed peers can therefore export identical bytes when every emitted
message was delivered and accepted. A dropped message can produce different
local transcripts. Local timeout and failure observations remain outside the
shared-message transcript.

The transcript is not signed. It is a canonical record of one participant's
accepted bytes, not consensus, identity proof, non-repudiation, proof of
delivery, or proof of motive.

\section{Conditional Correctness}

\textbf{Proposition.} Suppose Alice and Bob are honest, begin with identical
valid public context, successfully compute their commitments, and eventually
receive every canonical message before either local timeout. Then both enter
\texttt{COMPLETE}, accept the same pair \((m_A,m_B)\), and report the same
outcome.

\textbf{Proof.} Each party's move and nonce are fixed during initialization.
Both commitments are therefore available before either honest reveal is
emitted. Each received reveal is accepted only if recomputation in the agreed
context equals the corresponding earlier digest. Hence Alice and Bob accept
the same two encoded moves. Both evaluate the same deterministic expression
\(d=(m_A-m_B+3)\bmod 3\). For the nine possible pairs, \(d=0\) exactly for
equal moves; \(d=1\) for paper over rock, scissors over paper, and rock over
scissors; and \(d=2\) for the reverse cases. Therefore both report the same
tie or winner. \(\square\)

This proposition is conditional. It says nothing about completion under
message loss, malicious abort, compromised randomness, or a broken hash
assumption.

\section{Claims and Evidence Status}

\begin{center}
\small
\begin{tabular}{|>{\raggedright\arraybackslash}p{0.20\textwidth}|
                  >{\raggedright\arraybackslash}p{0.27\textwidth}|
                  >{\raggedright\arraybackslash}p{0.39\textwidth}|}
\hline
Property & Status & Boundary \\
\hline
Honest correctness & Conditional proposition & Requires identical context, honest parties, successful operations, and eventual delivery before timeout. \\
\hline
Hiding before reveal & Computational security argument & Depends on a fresh secret 256-bit nonce and a random-oracle-style assumption; timing and context remain visible. \\
\hline
Binding after commit & Computational security argument & Depends on collision and alternate-opening difficulty for the exact canonical encoding; no formal reduction is supplied. \\
\hline
Knowledge of opening & Not achieved at commit time & A sender may submit a random digest and later fail to open it. \\
\hline
Context separation & Protocol invariant & Domain, versions, session, round, roles, and payload type are bound into the digest; this is not authentication. \\
\hline
Parsing and state behavior & Implementation evidence only & Unit and adversarial tests can establish enumerated behavior, not cryptographic hardness. \\
\hline
Strong fairness and guaranteed output & Not achieved & The second revealer may learn the result and then withhold its reveal. \\
\hline
Transcript accountability & Not achieved & The local unsigned transcript does not prove identity, delivery, or motive to a third party. \\
\hline
\end{tabular}
\end{center}

\section{Security-Argument Status}

\subsection{Hiding argument}

The move alone has only three possibilities. Hiding therefore comes from the
fresh secret nonce, not from the move and not from the word ``hash.'' In a
random-oracle-style model, an observer trying to distinguish the committed
move must query the correct context, move, and 256-bit nonce or obtain an
equivalent collision. With an independently sampled nonce that remains secret,
this search is intended to be infeasible for a classical computationally
bounded observer.

This is only a heuristic security argument for a concrete hash function. The
guarantee is not information theoretic, and the argument does not formally
prove that BLAKE2b behaves as a random oracle. Nonce reuse, prediction,
logging, early disclosure, or endpoint compromise invalidates the argument.

\subsection{Binding argument}

To open one digest as two different canonical messages, a committer would need
two distinct encodings that produce the same 32-byte digest. For a fixed
honest commitment, an adversary seeking another opening faces an
alternate-opening or second-preimage problem. For a malicious committer that
chooses both encodings, the relevant concern is collision finding.

The encoding removes simple ambiguity attacks, but it does not prove the
assumed properties of the hash. No reduction from a formal commitment game to
a named BLAKE2b assumption is supplied here.

\subsection{Context separation}

Changing a protocol version, session, round, role, or payload type changes the
canonical preimage. This prevents an otherwise valid opening from being
accepted in a different context when the parser and state machine enforce the
same fields. It does not stop an unauthenticated attacker from inventing a new
message; that remains the transport layer's responsibility.

\section{Attacks and Counterexamples}

The implementation and chapter explicitly preserve the following failures:

\begin{itemize}
    \item \textbf{Commit without knowledge.} A malicious peer can send a
    context-valid random digest for which it knows no opening.
    \item \textbf{Wrong opening.} Changing the move or nonce causes
    verification failure, but does not compel a correct opening.
    \item \textbf{Cross-context replay.} A message from another session,
    version, role, or round must be rejected; context fields are not identity
    credentials.
    \item \textbf{Premature reveal and duplicates.} Version 1 enters a failed
    terminal state for an established in-session violation rather than
    recovering.
    \item \textbf{Silence.} A timeout is only the local observation that a
    message did not arrive. It proves neither cause nor intent.
    \item \textbf{Same-process demonstration.} Two independent state objects
    help test ownership and chronology, but one teaching process is not a real
    confidentiality boundary.
    \item \textbf{Selective abort.} The second revealer can learn the result
    and refuse to complete. This is developed next.
\end{itemize}

\section{Selective Abort, Timeout, and Fairness}

Consider the following valid prefix of an execution:

\begin{enumerate}
    \item Alice and Bob accept both commitments.
    \item Alice reveals first.
    \item Bob validates Alice's opening. Because Bob already knows his own
    move, Bob can now compute the result.
    \item Bob's participant-local API still refuses to report a completed
    outcome because Bob has not emitted his reveal. This API rule does not
    erase what the malicious host can compute from its own memory.
    \item If the result is unfavorable, Bob withholds his reveal or sends an
    abort message.
    \item Alice has only Bob's hiding commitment. She records the abort or
    times out without learning Bob's move or the result. Neither local session
    reaches \texttt{COMPLETE}.
\end{enumerate}

Bob has learned more than Alice and can bias which games reach completion. An
explicit abort makes the peer's declaration visible within the assumed
authenticated session, but it does not deliver the missing move. Silence gives
Alice only a local timeout. The unsigned transcript proves neither delivery nor
motive to a third party.

Declaring the incomplete game a tie or automatically awarding a winner would
not be a cryptographic repair. It would add an enforcement and utility
mechanism whose authority, incentives, and abuse cases require a separate
model.

\section{C Implementation Companion}

The companion is divided so that the game is not hidden inside a library call:

\begin{description}
    \item[Commitment interface.] \leavevmode\\
    \path{include/ac/commitment.h} exposes typed
    context, nonce, digest, computation, verification, and clearing operations.
    \item[Commitment implementation.] \leavevmode\\
    \path{src/primitives/commitment.c}
    performs canonical incremental hashing and checked library operations.
    \item[RPS interface.] \leavevmode\\
    \path{include/ac/rps.h} exposes participant-local
    initialization, message creation and receipt, timeout, typed abort reason,
    outcome, transcript, and cleanup operations.
    \item[RPS implementation.] \leavevmode\\
    \path{src/protocols/rps.c} owns the parser,
    state transitions, opening verification, outcome, and transcript.
    \item[Teaching driver.] \leavevmode\\
    \path{src/two_oracles_play_rock_paper_scissors/rps_demo.c} creates two
    separate local states and transfers serialized messages between them. It is
    not a network or process-isolation simulation.
    \item[Tests.] The test artifacts are separated by purpose:
    \begin{itemize}
        \item \path{tests/test_commitment.c}\par
        \noindent Deterministic commitment and mutation checks.\par
        \item \path{tests/test_rps.c}\par
        \noindent Honest games and core state behavior.\par
        \item \path{tests/test_rps_adversarial.c}\par
        Malformed, replay, invalid-opening, abort, timeout, and transcript
        attacks.\par
        \item \path{test-vectors/commitment-v1.md}\par
        \noindent The frozen encoding vector.\par
    \end{itemize}
\end{description}

The public interface maps to the protocol as follows:

\begin{itemize}
    \item initialize with \path{ac_rps_session_init}; deterministic tests alone
    use \path{ac_rps_session_init_with_nonce_for_test};
    \item emit messages with \path{ac_rps_make_commit},
    \path{ac_rps_make_reveal}, and \path{ac_rps_make_abort};
    \item process every peer message through \path{ac_rps_receive};
    \item observe timeout, abort reason, and outcome through
    \path{ac_rps_mark_timeout}, \path{ac_rps_get_abort_reason}, and
    \path{ac_rps_get_outcome}; and
    \item export or clear local state through \path{ac_rps_export_transcript}
    and \path{ac_rps_session_clear}.
\end{itemize}

Ordinary execution obtains nonces from libsodium's operating-system-backed
random interface. Tests inject fixed nonces through a separately named path.
The library backend supplies hashing, random bytes, constant-time digest
comparison, and memory-clearing helpers; the protocol itself remains visible
as C11 state and byte encoding.

The demonstration must print an educational, non-production warning, state the
authenticated-transport assumption, and show both honest completion and the
asymmetric selective-abort trace. It must not print a move or nonce before the
corresponding reveal.

\section{Tests and Experiments}

The companion test suite is intended to check:

\begin{itemize}
    \item a deterministic commitment vector with the full canonical preimage;
    \item mutation of every commitment context field, payload, nonce, and
    digest;
    \item all nine honest move pairs from both perspectives;
    \item Alice-first and Bob-first legal message orders;
    \item truncation, trailing data, invalid values, and length mismatches;
    \item duplicate, replayed, reflected, reordered, and premature messages;
    \item random commitments with no known opening and invalid openings;
    \item abort before commitment, after both commitments, and after one reveal;
    \item timeout, dropped delivery, terminal-state behavior, and transcript
    disagreement; and
    \item the selective-abort asymmetry described above.
\end{itemize}

A passing test establishes behavior for the exercised implementation and
configuration. It does not establish computational hiding, binding, fairness,
or production readiness. Tool versions and actual pass/fail results belong in
the implementation report rather than being predicted by this chapter.

\section{Ethical and Institutional Analysis}

The protocol replaces one narrow trust relation: neither party must trust the
other not to exploit an early plaintext move. It does so by fixing chronology
in a transcript that can later reject inconsistent openings.

It also creates choices that cryptography cannot settle. Who defines the
timeout? What consequence follows from refusal? Can a weaker participant
leave without punishment? Who can audit an alleged abort? Does an external
penalty authority have too much power? A technically correct commitment does
not answer those questions.

The selective-abort trace is especially important institutionally. Detecting
that cooperation stopped is not the same as restoring symmetry, and attaching
punishment may create a new coercive institution. Later mechanisms must be
judged both for their cryptographic properties and for who controls them.

\section{What Cryptography Does Not Solve}

Cryptography can help fix a move before disclosure and detect an inconsistent
opening. It does not:

\begin{itemize}
    \item decide which move is good;
    \item make either participant choose randomly or rationally;
    \item force a participant or network to deliver a reveal;
    \item give both parties the result at the same instant;
    \item establish human identity or a trustworthy channel;
    \item explain why a participant aborted;
    \item prevent learning from the public history of completed games; or
    \item turn a game rule or enforcement mechanism into a just one.
\end{itemize}

The protocol can constrain a small universe. It cannot choose the purpose of
that universe for its inhabitants.

\section{Limitations}

This is a two-party, one-round educational construction. The security
discussion is classical and non-composable. The exact hash commitment has no
formal proof here. The channel assumption is not implemented. The state
machine fails closed rather than recovering from arbitrary reordering or
duplication. Local transcripts may disagree under message loss. The teaching
driver runs both participant states in one process. The protocol provides no
fair output, guaranteed delivery, crash recovery, persistence, concurrency,
coercion resistance, or strategy privacy after reveal.

These limitations are part of the result. Removing them requires additional
mechanisms and assumptions, not stronger adjectives.

\section{Open Problems}

\begin{enumerate}
    \item Does the original oracle sentence intend a problem beyond the narrow
    timing interpretation used here?
    \item Can the exact commitment be replaced by a comparably clear
    construction with a directly applicable formal proof and equally minimal
    implementation dependency?
    \item What is the smallest mechanism that improves fairness without
    introducing a trusted adjudicator, coercive penalty, or hidden new
    assumption?
    \item How should an authenticated transport expose delay, duplication, and
    retry without making the protocol state ambiguous?
    \item What information about strategy is inevitably leaked by repeated
    completed rounds?
    \item Under what institutional conditions is an abort consequence
    legitimate rather than coercive?
\end{enumerate}

\section{Exercises}

\begin{enumerate}
    \item Enumerate all nine move pairs and verify the outcome formula.
    \item Show why hashing only one of three move bytes does not hide the move.
    \item Construct two messages from different sessions and identify every
    field that prevents cross-session acceptance.
    \item Write the local states observed by Alice and Bob in the selective-
    abort trace. Which facts are shared, and which are merely local?
    \item Explain why refusing to return an outcome from an API before both
    reveals does not prevent a malicious host from learning it early.
    \item Propose one fairness mechanism. List the new trusted power,
    assumption, cost, and abuse case that it introduces.
\end{enumerate}

The bibliography for this chapter is emitted by the book or standalone-section
wrapper. The chapter intentionally contains no independent bibliography
declaration.
